
\documentclass{article}
\usepackage{colortbl}
\usepackage{makecell}
\usepackage{multirow}
\usepackage{supertabular}

\begin{document}

\newcounter{utterance}

\centering \large Interaction Transcript for game `dond', experiment `coop\_en', episode 10 with Aya{-}Expanse{-}32B{-}t0.0.
\vspace{24pt}

{ \footnotesize  \setcounter{utterance}{1}
\setlength{\tabcolsep}{0pt}
\begin{supertabular}{c@{$\;$}|p{.15\linewidth}@{}p{.15\linewidth}p{.15\linewidth}p{.15\linewidth}p{.15\linewidth}p{.15\linewidth}}
    \# & $\;$A & \multicolumn{4}{c}{Game Master} & $\;\:$B\\
    \hline

    \theutterance \stepcounter{utterance}  
    & & \multicolumn{4}{p{0.6\linewidth}}{
        \cellcolor[rgb]{0.9,0.9,0.9}{
            \makecell[{{p{\linewidth}}}]{
                \texttt{\tiny{[P1$\langle$GM]}}
                \texttt{You are playing a cooperative negotiation game in which you have to agree on how to divide a set of items among you and another player.} \\
\\ 
\texttt{Rules:} \\
\texttt{(a) You and the other player are given a set of items. Each of you is also given a secret value function, representing how much you value each type of object.} \\
\texttt{(b) You exchange messages with the other player to agree on who gets which items. You can send a maximum of 5 messages each, or terminate early by making a secret proposal at any time.} \\
\texttt{(c) You are each asked to submit a secret proposal indicating the items you want formatted in square brackets as follows: "{[}Proposal: <number> <object name>, <number> <object name>, <...>{]}"} \\
\texttt{(d) If your proposals are complementary, i.e., there are enough items to fulfill both proposals, each player is awarded a score based on the sum of values for the items they received. Otherwise, both of you get zero points.} \\
\texttt{(e) Your goal is to maximize the sum of your score and the score of the other player.} \\
\\ 
\texttt{Let us start.} \\
\\ 
\texttt{The set of available items is:} \\
\\ 
\texttt{1 glass, 2 wires, 2 chalks} \\
\\ 
\texttt{Your secret value function is:} \\
\\ 
\texttt{glass: 8, wire: 1, chalk: 0} \\
\\ 
\texttt{IMPORTANT: Your messages, unless it is the secret proposal, are directly transmitted to the other player, so do not include any response to the rules or text announcing your message. To make a secret proposal, use the indicated format. Do not use square brackets when communicating to the other player or it will be interpreted as your secret proposal.} \\
            }
        }
    }
    & \\ \\

    \theutterance \stepcounter{utterance}  
    & \multicolumn{4}{p{0.6\linewidth}}{
        \cellcolor[rgb]{0.9,0.9,0.9}{
            \makecell[{{p{\linewidth}}}]{
                \texttt{\tiny{[P1$\rangle$GM]}}
                \texttt{Understanding the value of items to myself, and aiming to maximize the total score, I'd like to propose a cooperative division. Given the limited communication, I'll suggest a fair split right away, assuming the other player has a similar goal.} \\
\\ 
\texttt{Message 1: I propose we give them 2 chalks and I take the glass and 2 wires.} \\
\\ 
\texttt{This division respects our value functions, as per my understanding. If they value chalks more, this should be a favorable proposition for both of us.} \\
\\ 
\texttt{After their response, I can adjust my proposal or make a secret proposal if necessary, but initially, I believe this is a good starting point.} \\
            }
        }
    }
    & & \\ \\

    \theutterance \stepcounter{utterance}  
    & & & \multicolumn{2}{p{0.3\linewidth}}{
        \cellcolor[rgb]{0.9,0.9,0.9}{
            \makecell[{{p{\linewidth}}}]{
                \texttt{\tiny{[GM$|$GM]}}
                \texttt{continue} \\
            }
        }
    }
    & & \\ \\

    \theutterance \stepcounter{utterance}  
    & & & \multicolumn{2}{p{0.3\linewidth}}{
        \cellcolor[rgb]{0.9,0.9,0.9}{
            \makecell[{{p{\linewidth}}}]{
                \texttt{\tiny{[GM$|$GM]}}
                \texttt{Understanding the value of items to myself, and aiming to maximize the total score, I'd like to propose a cooperative division. Given the limited communication, I'll suggest a fair split right away, assuming the other player has a similar goal.} \\
\\ 
\texttt{Message 1: I propose we give them 2 chalks and I take the glass and 2 wires.} \\
\\ 
\texttt{This division respects our value functions, as per my understanding. If they value chalks more, this should be a favorable proposition for both of us.} \\
\\ 
\texttt{After their response, I can adjust my proposal or make a secret proposal if necessary, but initially, I believe this is a good starting point.} \\
            }
        }
    }
    & & \\ \\

    \theutterance \stepcounter{utterance}  
    & & \multicolumn{4}{p{0.6\linewidth}}{
        \cellcolor[rgb]{0.9,0.9,0.9}{
            \makecell[{{p{\linewidth}}}]{
                \texttt{\tiny{[GM$\rangle$P2]}}
                \texttt{You are playing a cooperative negotiation game in which you have to agree on how to divide a set of items among you and another player.} \\
\\ 
\texttt{Rules:} \\
\texttt{(a) You and the other player are given a set of items. Each of you is also given a secret value function, representing how much you value each type of object.} \\
\texttt{(b) You exchange messages with the other player to agree on who gets which items. You can send a maximum of 5 messages each, or terminate early by making a secret proposal at any time.} \\
\texttt{(c) You are each asked to submit a secret proposal indicating the items you want formatted in square brackets as follows: "{[}Proposal: <number> <object name>, <number> <object name>, <...>{]}"} \\
\texttt{(d) If your proposals are complementary, i.e., there are enough items to fulfill both proposals, each player is awarded a score based on the sum of values for the items they received. Otherwise, both of you get zero points.} \\
\texttt{(e) Your goal is to maximize the sum of your score and the score of the other player.} \\
\\ 
\texttt{Let us start.} \\
\\ 
\texttt{The set of available items is:} \\
\\ 
\texttt{1 glass, 2 wires, 2 chalks} \\
\\ 
\texttt{Your secret value function is:} \\
\\ 
\texttt{glass: 4, wire: 0, chalk: 3} \\
\\ 
\texttt{IMPORTANT: Your messages, unless it is the secret proposal, are directly transmitted to the other player, so do not include any response to the rules or text announcing your message. To make a secret proposal, use the indicated format. Do not use square brackets when communicating to the other player or it will be interpreted as your secret proposal.} \\
\\ 
\\ 
\texttt{Understanding the value of items to myself, and aiming to maximize the total score, I'd like to propose a cooperative division. Given the limited communication, I'll suggest a fair split right away, assuming the other player has a similar goal.} \\
\\ 
\texttt{Message 1: I propose we give them 2 chalks and I take the glass and 2 wires.} \\
\\ 
\texttt{This division respects our value functions, as per my understanding. If they value chalks more, this should be a favorable proposition for both of us.} \\
\\ 
\texttt{After their response, I can adjust my proposal or make a secret proposal if necessary, but initially, I believe this is a good starting point.} \\
            }
        }
    }
    & \\ \\

    \theutterance \stepcounter{utterance}  
    & & & \multicolumn{4}{p{0.6\linewidth}}{
        \cellcolor[rgb]{0.9,0.9,0.9}{
            \makecell[{{p{\linewidth}}}]{
                \texttt{\tiny{[GM$\langle$P2]}}
                \texttt{Message 2: Interesting proposal. Before I respond, I want to clarify our value functions to ensure we're on the same page. Could you confirm if your valuation is: glass > chalk > wire?} \\
\\ 
\texttt{Assuming your valuation aligns with mine (glass being the most valuable, followed by chalk, and wire being the least valuable), I'd like to suggest a slight modification to ensure a fair split.} \\
\\ 
\texttt{Message 3: How about we split the chalks evenly, with me taking the glass and one wire, and you taking the other wire and one chalk? This way, we both get items we value highly.} \\
\\ 
\texttt{Proposing a division where both parties receive items they value can foster cooperation and maximize our combined score. I believe this split respects our value functions while ensuring a fair outcome.} \\
\\ 
\texttt{I will await their response and be prepared to make a secret proposal if necessary after their next message.} \\
            }
        }
    }
     \\ \\

    \theutterance \stepcounter{utterance}  
    & & & \multicolumn{2}{p{0.3\linewidth}}{
        \cellcolor[rgb]{0.9,0.9,0.9}{
            \makecell[{{p{\linewidth}}}]{
                \texttt{\tiny{[GM$|$GM]}}
                \texttt{continue} \\
            }
        }
    }
    & & \\ \\

    \theutterance \stepcounter{utterance}  
    & & & \multicolumn{2}{p{0.3\linewidth}}{
        \cellcolor[rgb]{0.9,0.9,0.9}{
            \makecell[{{p{\linewidth}}}]{
                \texttt{\tiny{[GM$|$GM]}}
                \texttt{Message 2: Interesting proposal. Before I respond, I want to clarify our value functions to ensure we're on the same page. Could you confirm if your valuation is: glass > chalk > wire?} \\
\\ 
\texttt{Assuming your valuation aligns with mine (glass being the most valuable, followed by chalk, and wire being the least valuable), I'd like to suggest a slight modification to ensure a fair split.} \\
\\ 
\texttt{Message 3: How about we split the chalks evenly, with me taking the glass and one wire, and you taking the other wire and one chalk? This way, we both get items we value highly.} \\
\\ 
\texttt{Proposing a division where both parties receive items they value can foster cooperation and maximize our combined score. I believe this split respects our value functions while ensuring a fair outcome.} \\
\\ 
\texttt{I will await their response and be prepared to make a secret proposal if necessary after their next message.} \\
            }
        }
    }
    & & \\ \\

    \theutterance \stepcounter{utterance}  
    & & \multicolumn{4}{p{0.6\linewidth}}{
        \cellcolor[rgb]{0.9,0.9,0.9}{
            \makecell[{{p{\linewidth}}}]{
                \texttt{\tiny{[P1$\langle$GM]}}
                \texttt{Message 2: Interesting proposal. Before I respond, I want to clarify our value functions to ensure we're on the same page. Could you confirm if your valuation is: glass > chalk > wire?} \\
\\ 
\texttt{Assuming your valuation aligns with mine (glass being the most valuable, followed by chalk, and wire being the least valuable), I'd like to suggest a slight modification to ensure a fair split.} \\
\\ 
\texttt{Message 3: How about we split the chalks evenly, with me taking the glass and one wire, and you taking the other wire and one chalk? This way, we both get items we value highly.} \\
\\ 
\texttt{Proposing a division where both parties receive items they value can foster cooperation and maximize our combined score. I believe this split respects our value functions while ensuring a fair outcome.} \\
\\ 
\texttt{I will await their response and be prepared to make a secret proposal if necessary after their next message.} \\
            }
        }
    }
    & \\ \\

    \theutterance \stepcounter{utterance}  
    & \multicolumn{4}{p{0.6\linewidth}}{
        \cellcolor[rgb]{0.9,0.9,0.9}{
            \makecell[{{p{\linewidth}}}]{
                \texttt{\tiny{[P1$\rangle$GM]}}
                \texttt{Given your clarification and proposal, let's summarize our understanding:} \\
\\ 
\texttt{{-} You value items in the order: glass > chalk > wire.} \\
\texttt{{-} I have the same valuation: glass > wire > chalk.} \\
\\ 
\texttt{Your proposed division is:} \\
\texttt{{-} You: {[}1 glass, 1 wire{]}} \\
\texttt{{-} Me: {[}1 wire, 1 chalk{]}} \\
\\ 
\texttt{This split acknowledges the higher value we both place on the glass and wires. By dividing the chalks evenly, we ensure neither of us is left with an item we value least.} \\
\\ 
\texttt{**Secret Proposal:**} \\
\texttt{Considering our shared understanding and the proposed division, I submit the following secret proposal:} \\
\\ 
\texttt{"Proposal: 1 glass, 1 wire"} \\
\\ 
\texttt{This proposal reflects the division we discussed and maximizes our combined score based on our value functions. I believe this is a fair and cooperative outcome.} \\
            }
        }
    }
    & & \\ \\

    \theutterance \stepcounter{utterance}  
    & & & \multicolumn{2}{p{0.3\linewidth}}{
        \cellcolor[rgb]{0.9,0.9,0.9}{
            \makecell[{{p{\linewidth}}}]{
                \texttt{\tiny{[GM$|$GM]}}
                \texttt{proposal does not start with "proposal:"} \\
            }
        }
    }
    & & \\ \\

\end{supertabular}
}

\end{document}
